\chapter{Бикомплекс полной степени и no-go единственной Hodge-метрики}
\label{ch:v7-bicomplex-total-degree-hodge-metric-gate}

\section{Проверяемое усиление}

Предыдущий гейт поместил рёберную и связывающую кривизны в один носитель
размерности $96$, но оставил семейство положительных метрик
\begin{equation}
 G_\eta=P_E+\eta P_L,
 \qquad \eta>0,
 \label{eq:v7-bicomplex-metric-family}
\end{equation}
где $P_E$ и $P_L$ являются ортогональными проекторами на блоки размерностей
$54$ и $42$. Теперь проверим, устраняют ли эту свободу структура
бикомплекса, полная степень, Real-сопряжение или Hodge-звезда.

Для двойного комплекса требуются
\begin{equation}
 d_h^2=d_v^2=0,
 \qquad d_hd_v+d_vd_h=0,
 \qquad D=d_h+d_v.
 \label{eq:v7-bicomplex-relations}
\end{equation}
Однако для любого $c>0$ замена
\begin{equation}
 d_v\longmapsto d_v^{(c)}=c\,d_v
 \label{eq:v7-bicomplex-vertical-rescaling}
\end{equation}
сохраняет все тождества \eqref{eq:v7-bicomplex-relations}.
Следовательно, сами алгебраические аксиомы бикомплекса не фиксируют
отношение длин горизонтального и вертикального направлений.

\section{Точный квадратный контроль}

На базисе бикомплекса
\begin{equation}
 (0,0),(1,0),(0,1),(1,1)
 \end{equation}
возьмём стрелки
\begin{equation}
 d_h e_{00}=e_{10},\quad d_h e_{01}=e_{11},\qquad
 d_v^{(c)}e_{00}=c e_{01},\quad d_v^{(c)}e_{10}=-c e_{11}.
 \label{eq:v7-bicomplex-square-control}
\end{equation}
Для $c=1/8,1/2,1,2,8$ машинный расчёт даёт нулевые с точностью вычислений
остатки нильпотентности и антикоммутатора. Если
\begin{equation}
 N_{\rm tot}=\operatorname{diag}(0,1,1,2),
 \end{equation}
то при каждом $c$
\begin{equation}
 [N_{\rm tot},D_c]=D_c.
 \label{eq:v7-bicomplex-total-degree}
\end{equation}
Таким образом, даже целочисленная полная степень не видит относительного
масштаба двух дифференциалов. Это согласуется с общим алгебраическим языком
двойных комплексов как сумм квадратов и зигзагов
\cite{v7-bicomplex-stelzig} и с параметризованными квадрат-нуль
дифференциалами в Hodge-теории \cite{v7-bicomplex-tardini-tomassini}.

\section{Hodge-звезда также допускает семейство}

Положительность, Real-совместимость и условие $\star^2=1$ сами по себе не
устраняют \eqref{eq:v7-bicomplex-metric-family}. На удвоенном пространстве
$\mathcal K\oplus\mathcal K^\vee$ положим
\begin{equation}
 \star_\eta=
 \begin{pmatrix}
 0&G_\eta^{-1}\\
 G_\eta&0
 \end{pmatrix},
 \qquad
 H_\eta=\operatorname{diag}(G_\eta,G_\eta^{-1}).
 \label{eq:v7-bicomplex-hodge-family}
\end{equation}
Тогда для каждого $\eta>0$
\begin{equation}
 \star_\eta^2=1,
 \qquad
 \star_\eta^*H_\eta\star_\eta=H_\eta.
 \label{eq:v7-bicomplex-hodge-isometry}
\end{equation}
Все матрицы вещественны и коммутируют с полной степенью и текущим
комплексным сопряжением. Машинные остатки равны нулю для
$\eta=1/8,1/2,1,2,8$. Значит, Hodge-звезда становится единственной только
после независимого задания геометрической метрики; она не выводит эту
метрику из bigrading. В стандартной литературе это разделение также явно:
сначала задаётся двойной комплекс, а Hodge-операторы строятся после выбора
эрмитовой метрики \cite{v7-bicomplex-tardini-tomassini,
v7-bicomplex-popovici-stelzig-ugarte}.

\section{Размерностный дефект}

Обменная изометрия между текущими компонентами невозможна, поскольку
\begin{equation}
 \rank P_E=54,
 \qquad \rank P_L=42.
 \end{equation}
Максимальная изометрическая инъекция
$U_L:\mathbb C^{42}\to\mathbb C^{54}$ удовлетворяет
\begin{equation}
 U_L^*U_L=I_{42},
 \qquad
 \rank(I_{54}-U_LU_L^*)=12.
 \label{eq:v7-bicomplex-defect-twelve}
\end{equation}
Поэтому никакая унитарная Hodge-операция не может просто поменять два блока
местами. Неприводимость текущего носителя также отсутствует: $P_E$ и $P_L$
коммутируют со всеми уже построенными блок-диагональными операторами.

Число $12$ заслуживает отдельной проверки:
\begin{equation}
 54-42=12=\dim_{\mathbb C}E_{\rm aff}.
 \label{eq:v7-bicomplex-affine-dimension-clue}
\end{equation}
Это пока только размерностное совпадение. Ни отображение
$E_{\rm aff}\to\ker U_L^*$, ни совместимость бимодульных меток, ни новый
смешанный дифференциал ещё не выведены.

\section{Физическое следствие}

На полном 27-мерном гессиане веса $\eta=1/4,1,4$ сохраняют
\begin{equation}
 (n_-,n_0,n_+)_{0}=(7,0,20),
 \qquad
 (n_-,n_0,n_+)_{*}=(0,0,27).
\end{equation}
Но вакуумные спектры не связаны одним общим масштабом. После наилучшего
глобального перескалирования относительная невязка спектров при
$\eta=1/4$ и $\eta=4$ равна
\begin{equation}
 0.3181422930\ldots.
 \label{eq:v7-bicomplex-nonhomothetic-spectrum}
\end{equation}
Следовательно, качественный вакуум остаётся строгим, а отношения масс
действительно зависят от невыведенной метрики.

\begin{theorem}[No-go полной степени]
Алгебраические аксиомы бикомплекса, оператор полной степени, текущее
Real-сопряжение и существование изометрической Hodge-инволюции совместимы с
непрерывным семейством положительных относительных метрик $G_\eta$.
Следовательно, они не фиксируют $\eta$ и не выводят отношения масс.
\end{theorem}

\begin{remark}
Отрицательный результат не переоткрывает качественную проблему. При всех
$\eta\ge0$ сохраняются правильный запуск и устойчивый вакуум. Он запрещает
только превращать выбранное значение $\eta$ в количественное предсказание.
\end{remark}

\begin{nogocriterion}[Запрет метрики из одной bigrading]
Нельзя объявлять $\eta=1$ следствием total-degree или Hodge-звезды без
оператора, смешивающего два бистепенных блока. Любой такой оператор обязан
быть выведен из существующего физического модуля, пройти Real-условия и не
внести новые произвольные массы.
\end{nogocriterion}

\begin{protocol}
\begin{enumerate}
 \item квадратный бикомплекс при пяти масштабах: \textbf{прошёл};
 \item $D_c^2=0$: \textbf{для всех $c$};
 \item $[N_{\rm tot},D_c]=D_c$: \textbf{для всех $c$};
 \item семейство положительных метрик $G_\eta$: \textbf{существует};
 \item Real-совместимая Hodge-инволюция: \textbf{для всех $\eta>0$};
 \item единственность относительной метрики: \textbf{не получена};
 \item унитарный обмен $54\leftrightarrow42$: \textbf{невозможен};
 \item размерностный дефект: \textbf{$12$};
 \item совпадение с $\dim E_{\rm aff}$: \textbf{получено};
 \item каноническое аффинное отождествление: \textbf{не получено};
 \item нулевая сигнатура: \textbf{$(7,0,20)$};
 \item вакуумная сигнатура: \textbf{$(0,0,27)$};
 \item массовые отношения: \textbf{не выведены};
 \item следующий гейт: \textbf{аффинное дефектное дополнение бикомплекса}.
\end{enumerate}
\end{protocol}

Машинный сертификат сохранён в
\path{s2t_v7_bicomplex_total_degree_hodge_metric_gate_results.json}.

\begin{thebibliography}{9}
\bibitem{v7-bicomplex-stelzig}
J.~Stelzig,
\emph{On the Structure of Double Complexes},
J. London Math. Soc. 104 (2021), 956--988; arXiv:1812.00865.
\bibitem{v7-bicomplex-tardini-tomassini}
N.~Tardini, A.~Tomassini,
\emph{Differential operators on almost-Hermitian manifolds and harmonic
forms}, Complex Manifolds 7 (2020), 106--128; arXiv:1909.06569.
\bibitem{v7-bicomplex-popovici-stelzig-ugarte}
D.~Popovici, J.~Stelzig, L.~Ugarte,
\emph{Higher-Page Bott--Chern and Aeppli Cohomologies and Applications},
J. Reine Angew. Math. 777 (2021), 157--194; arXiv:2007.03320.
\end{thebibliography}